% Owner: TV3 (P5.3b) -- Reviewer: TV1
\subsection{Rollout Policies and Simulation Improvements}
\label{subsec:rollouts}

A rollout plays from the new node to the end of the game, or to a depth limit, and returns a value in $[-1, 1]$ for the player to move at its start. We compare five policies of increasing knowledge:
\begin{itemize}
    \item \textbf{L0}: uniformly random legal moves until the game ends.
    \item \textbf{L1}: if the player to move can win at once, it does so and the rollout ends; otherwise, if the opponent threatens to win at once, it blocks; otherwise it plays a random legal move.
    \item \textbf{L2}: like L1, but the random moves are drawn only from the candidate cells near existing stones.
    \item \textbf{L3}: like L2, but the remaining move is sampled from a softmax over the local pattern score, $P(m \mid s) \propto \exp(h(s,m)/T)$ with $T = 1$.
    \item \textbf{L4}: L3 truncated after $d_{\mathrm{cut}} = 8$ moves; the position is then scored by the evaluation function of Section~\ref{subsec:alphabeta} as $\tanh(\mathrm{eval}(s)/\tau)$ with $\tau = 2000$, so that a value replaces the rest of the random game.
\end{itemize}
The Gomoku rollouts work on a flat array with a two-cell border instead of the game object, and they track the endpoints of threats incrementally, so that each move checks at most eight cells instead of rescanning the board.

\textbf{Simulation improvements (bonus B4).} The \emph{extended forced moves} (L2-B4, L3-B4, L4-B4) add two tactical levels after win and block: a player that can make an open four or a double four wins the rollout at once, since such a threat cannot be stopped, and a player that faces an opponent's open three blocks it. The rollout also ends as a loss for the player to move when the opponent has two separate winning cells. \emph{Root parallelization} runs $K$ independent trees from the root in separate processes with different seeds and sums the visit counts and rewards of each root child before choosing the move.
